% GENERATED FILE. Edit canonical lessons, then run npm run generate:books.
% canonical-source-hash: fnv1a64:9189534e758c42d3
% canonical-lessons: SA-C293-dosah, SA-C293-dandayati, SA-C293-mocayati, SA-C293-karagaram, SA-C293-srnkhala, SA-R293-review-286-293

\chapter{The Case in Court}
\label{ch:sa-the-case-in-court}
\hlchaptermodality{293}

\begin{chapteropening}
\textbf{By the end of this chapter:} \emph{I can tell how a case ended and say what happens if someone is guilty.}

Complete the last lesson of chapter 293: I can tell how a case ended and say what happens if someone is guilty.
\end{chapteropening}

\section[\texorpdfstring{\sk{दोषः} (doṣaḥ)}{दोषः (doṣaḥ)}]{\sk{दोषः} (doṣaḥ) — guilt; a fault}

\label{lesson:SA-C293-dosah}



\begin{quote}
Before the new one: say the Sanskrit for a secret, then the Sanskrit for reveals.
\end{quote}

\subsection*{You'll want to know: \sk{दोषः}}
\textbf{\sk{दोषः}} — \emph{doṣaḥ} — \textquotedblleft{}guilt; a fault\textquotedblright{}.

\begin{quote}
\sk{चौरस्य} \sk{दोषः} \sk{स्पष्टः} \sk{अस्ति।}
\end{quote}

\begin{quote}
\emph{caurasya doṣaḥ spaṣṭaḥ asti.}
\end{quote}

\begin{quote}
The thief's guilt is clear.
\end{quote}

A masculine noun.

\begin{cousinweb}
\textbf{\sk{दोषः}} (\emph{doṣaḥ}) comes from the root √duṣ, \textquotedblleft{}to become bad\textquotedblright{}; MW gives \emph{doṣa} as a fault or defect and also as an offence or guilt (MW).
\end{cousinweb}

\begin{grammarlens}[title={What you've built}]
One more word for telling what happened.
\end{grammarlens}

\subsection*{Your turn}
\begin{itemize}
\raggedright
  \item \emph{Say it:} \emph{doṣaḥ}
  \item \emph{Say it:} \emph{doṣaḥ}, once more
  \item \emph{Say it:} \emph{caurasya doṣaḥ spaṣṭaḥ asti.} — The thief's guilt is clear.
  \item \emph{Recall:} say the Sanskrit for a secret, then the Sanskrit for reveals, then say \emph{doṣaḥ} again
\end{itemize}

\begin{tcolorbox}[breakable,colback=teal!4,colframe=teal!35!black,title={Before you move on}]
What does \sk{दोषः} mean? (\textquotedblleft{}Guilt; a fault\textquotedblright{}.) Say it once more.
\end{tcolorbox}
\section[\texorpdfstring{\sk{दण्डयति} (daṇḍayati)}{दण्डयति (daṇḍayati)}]{\sk{दण्डयति} (daṇḍayati) — punishes}

\label{lesson:SA-C293-dandayati}



\begin{quote}
Before the new one: say the Sanskrit for reveals, then the Sanskrit for guilt.
\end{quote}

\subsection*{You'll want to know: \sk{दण्डयति}}
\textbf{\sk{दण्डयति}} — \emph{daṇḍayati} — \textquotedblleft{}punishes\textquotedblright{}.

\begin{quote}
\sk{यदि} \sk{दोषः} \sk{अस्ति}, \sk{तर्हि} \sk{न्यायाधीशः} \sk{दण्डयति।}
\end{quote}

\begin{quote}
\emph{yadi doṣaḥ asti, tarhi nyāyādhīśaḥ daṇḍayati.}
\end{quote}

\begin{quote}
If there is guilt, then the judge punishes.
\end{quote}

Past (he …): \sk{सः} \sk{चौरं} \sk{दण्डितवान्}, \emph{saḥ cauraṃ daṇḍitavān}.

\begin{cousinweb}
\textbf{\sk{दण्डयति}} (\emph{daṇḍayati}) is built on \emph{daṇḍa}, \textquotedblleft{}a stick\textquotedblright{}, and so \textquotedblleft{}punishment\textquotedblright{}; MW gives \emph{daṇḍayati} as a verb made from it, to chastise or punish (MW).
\end{cousinweb}

\begin{grammarlens}[title={What you've built}]
One more word for telling what happened.
\end{grammarlens}

\subsection*{Your turn}
\begin{itemize}
\raggedright
  \item \emph{Say it:} \emph{daṇḍayati}
  \item \emph{Say it:} \emph{daṇḍayati}, once more
  \item \emph{Say it:} \emph{yadi doṣaḥ asti, tarhi nyāyādhīśaḥ daṇḍayati.} — If there is guilt, then the judge punishes.
  \item \emph{Recall:} say the Sanskrit for reveals, then the Sanskrit for guilt, then say \emph{daṇḍayati} again
\end{itemize}

\begin{tcolorbox}[breakable,colback=teal!4,colframe=teal!35!black,title={Before you move on}]
What does \sk{दण्डयति} mean? (\textquotedblleft{}Punishes\textquotedblright{}.) Say it once more.
\end{tcolorbox}
\section[\texorpdfstring{\sk{मोचयति} (mocayati)}{मोचयति (mocayati)}]{\sk{मोचयति} (mocayati) — sets free}

\label{lesson:SA-C293-mocayati}



\begin{quote}
Before the new one: say the Sanskrit for guilt, then the Sanskrit for punishes.
\end{quote}

\subsection*{You'll want to know: \sk{मोचयति}}
\textbf{\sk{मोचयति}} — \emph{mocayati} — \textquotedblleft{}sets free\textquotedblright{}.

\begin{quote}
\sk{न्यायाधीशः} \sk{बालं} \sk{मोचयति।}
\end{quote}

\begin{quote}
\emph{nyāyādhīśaḥ bālaṃ mocayati.}
\end{quote}

\begin{quote}
The judge sets the boy free.
\end{quote}

Past (he …): \sk{सः} \sk{बालं} \sk{मोचितवान्}, \emph{saḥ bālaṃ mocitavān}.

\begin{cousinweb}
\textbf{\sk{मोचयति}} (\emph{mocayati}) comes from the root √muc, \textquotedblleft{}to loose, release\textquotedblright{}, in its causative form, to set free (MW).
\end{cousinweb}

\begin{grammarlens}[title={What you've built}]
One more word for telling what happened.
\end{grammarlens}

\subsection*{Your turn}
\begin{itemize}
\raggedright
  \item \emph{Say it:} \emph{mocayati}
  \item \emph{Say it:} \emph{mocayati}, once more
  \item \emph{Say it:} \emph{nyāyādhīśaḥ bālaṃ mocayati.} — The judge sets the boy free.
  \item \emph{Recall:} say the Sanskrit for guilt, then the Sanskrit for punishes, then say \emph{mocayati} again
\end{itemize}

\begin{tcolorbox}[breakable,colback=teal!4,colframe=teal!35!black,title={Before you move on}]
What does \sk{मोचयति} mean? (\textquotedblleft{}Sets free\textquotedblright{}.) Say it once more.
\end{tcolorbox}
\section[\texorpdfstring{\sk{कारागारम्} (kārāgāram)}{कारागारम् (kārāgāram)}]{\sk{कारागारम्} (kārāgāram) — a prison}

\label{lesson:SA-C293-karagaram}



\begin{quote}
Before the new one: say the Sanskrit for punishes, then the Sanskrit for sets free.
\end{quote}

\subsection*{You'll want to know: \sk{कारागारम्}}
\textbf{\sk{कारागारम्}} — \emph{kārāgāram} — \textquotedblleft{}a prison\textquotedblright{}.

\begin{quote}
\sk{चौरः} \sk{कारागारे} \sk{वसति।}
\end{quote}

\begin{quote}
\emph{cauraḥ kārāgāre vasati.}
\end{quote}

\begin{quote}
The thief lives in prison.
\end{quote}

A neuter noun.

\begin{cousinweb}
\textbf{\sk{कारागारम्}} (\emph{kārāgāram}) is built from \emph{kārā}, \textquotedblleft{}confinement\textquotedblright{}, and \emph{āgāra}, \textquotedblleft{}a house\textquotedblright{}; MW gives \emph{kārāgāra} as a gaol (MW).
\end{cousinweb}

\begin{grammarlens}[title={What you've built}]
One more word for telling what happened.
\end{grammarlens}

\subsection*{Your turn}
\begin{itemize}
\raggedright
  \item \emph{Say it:} \emph{kārāgāram}
  \item \emph{Say it:} \emph{kārāgāram}, once more
  \item \emph{Say it:} \emph{cauraḥ kārāgāre vasati.} — The thief lives in prison.
  \item \emph{Recall:} say the Sanskrit for punishes, then the Sanskrit for sets free, then say \emph{kārāgāram} again
\end{itemize}

\begin{tcolorbox}[breakable,colback=teal!4,colframe=teal!35!black,title={Before you move on}]
What does \sk{कारागारम्} mean? (\textquotedblleft{}A prison\textquotedblright{}.) Say it once more.
\end{tcolorbox}
\section[\texorpdfstring{\sk{शृंखला} (śṛṅkhalā)}{शृंखला (śṛṅkhalā)}]{\sk{शृंखला} (śṛṅkhalā) — a chain}

\label{lesson:SA-C293-srnkhala}



\begin{quote}
Before the new one: say the Sanskrit for sets free, then the Sanskrit for a prison.
\end{quote}

\subsection*{You'll want to know: \sk{शृंखला}}
\textbf{\sk{शृंखला}} — \emph{śṛṅkhalā} — \textquotedblleft{}a chain\textquotedblright{}.

\begin{quote}
\sk{चौरस्य} \sk{पादे} \sk{शृंखला} \sk{अस्ति।}
\end{quote}

\begin{quote}
\emph{caurasya pāde śṛṅkhalā asti.}
\end{quote}

\begin{quote}
There is a chain on the thief's foot.
\end{quote}

A feminine noun.

\begin{cousinweb}
\textbf{\sk{शृंखला}} (\emph{śṛṅkhalā}) MW gives \emph{śṛṅkhalā} as a chain or fetter, without a derivation, so its origin is uncertain (MW).
\end{cousinweb}

\begin{grammarlens}[title={What you've built}]
One more word for telling what happened.
\end{grammarlens}

\subsection*{Your turn}
\begin{itemize}
\raggedright
  \item \emph{Say it:} \emph{śṛṅkhalā}
  \item \emph{Say it:} \emph{śṛṅkhalā}, once more
  \item \emph{Say it:} \emph{caurasya pāde śṛṅkhalā asti.} — There is a chain on the thief's foot.
  \item \emph{Recall:} say the Sanskrit for sets free, then the Sanskrit for a prison, then say \emph{śṛṅkhalā} again
\end{itemize}

\begin{tcolorbox}[breakable,colback=teal!4,colframe=teal!35!black,title={Before you move on}]
What does \sk{शृंखला} mean? (\textquotedblleft{}A chain\textquotedblright{}.) Say it once more.
\end{tcolorbox}
\section[(dialogue)]{Looking back: from A Theft in the Village to The Case in Court}

\label{lesson:SA-R293-review-286-293}



\begin{quote}
Say the Sanskrit for a theft and a night watchman.
\end{quote}

\subsection*{Your turn}
Say these aloud:

\begin{itemize}
\raggedright
  \item a theft, a night watchman, makes off with, a witness, runs away
  \item runs after, speed, creeps up to, stays awake, an uproar
  \item an offence, a suspicion, a charge in law, threatens, a weapon
  \item recognizes, describes, observes closely, a trace, before one's eyes
  \item a robber, plunders, attacks, a piece of jewellery, by force
  \item a rogue, cheats, a bribe, fraud, fake
  \item a spy, proof, examines, a secret, reveals
  \item guilt, punishes, sets free, a prison, a chain
\end{itemize}

\begin{tcolorbox}[breakable,colback=teal!4,colframe=teal!35!black,title={Before you move on}]
A theft? (\textbf{\sk{चौर्यम्}}, \emph{cauryam}.) A night watchman? (\textbf{\sk{यामिकः}}, \emph{yāmikaḥ}.) Makes off with? (\textbf{\sk{अपहरति}}, \emph{apaharati}.) A witness? (\textbf{\sk{साक्षी}}, \emph{sākṣī}.) Runs away? (\textbf{\sk{पलायते}}, \emph{palāyate}.) Runs after? (\textbf{\sk{अनुधावति}}, \emph{anudhāvati}.) Speed? (\textbf{\sk{वेगः}}, \emph{vegaḥ}.) Creeps up to? (\textbf{\sk{उपसर्पति}}, \emph{upasarpati}.) Stays awake? (\textbf{\sk{जागर्ति}}, \emph{jāgarti}.) An uproar? (\textbf{\sk{कोलाहलः}}, \emph{kolāhalaḥ}.) An offence? (\textbf{\sk{अपराधः}}, \emph{aparādhaḥ}.) A suspicion? (\textbf{\sk{सन्देहः}}, \emph{sandehaḥ}.) A charge in law? (\textbf{\sk{अभियोगः}}, \emph{abhiyogaḥ}.) Threatens? (\textbf{\sk{तर्जयति}}, \emph{tarjayati}.) A weapon? (\textbf{\sk{शस्त्रम्}}, \emph{śastram}.) Recognizes? (\textbf{\sk{अभिजानाति}}, \emph{abhijānāti}.) Describes? (\textbf{\sk{वर्णयति}}, \emph{varṇayati}.) Observes closely? (\textbf{\sk{निरीक्षते}}, \emph{nirīkṣate}.) A trace? (\textbf{\sk{चिह्नम्}}, \emph{cihnam}.) Before one's eyes? (\textbf{\sk{प्रत्यक्षम्}}, \emph{pratyakṣam}.) A robber? (\textbf{\sk{दस्युः}}, \emph{dasyuḥ}.) Plunders? (\textbf{\sk{लुण्ठति}}, \emph{luṇṭhati}.) Attacks? (\textbf{\sk{आक्रमति}}, \emph{ākramati}.) A piece of jewellery? (\textbf{\sk{आभरणम्}}, \emph{ābharaṇam}.) By force? (\textbf{\sk{बलात्}}, \emph{balāt}.) A rogue? (\textbf{\sk{धूर्तः}}, \emph{dhūrtaḥ}.) Cheats? (\textbf{\sk{वञ्चयति}}, \emph{vañcayati}.) A bribe? (\textbf{\sk{उत्कोचः}}, \emph{utkocaḥ}.) Fraud? (\textbf{\sk{कपटम्}}, \emph{kapaṭam}.) Fake? (\textbf{\sk{कृत्रिमः}}, \emph{kṛtrimaḥ}.) A spy? (\textbf{\sk{चरः}}, \emph{caraḥ}.) Proof? (\textbf{\sk{प्रमाणम्}}, \emph{pramāṇam}.) Examines? (\textbf{\sk{परीक्षते}}, \emph{parīkṣate}.) A secret? (\textbf{\sk{रहस्यम्}}, \emph{rahasyam}.) Reveals? (\textbf{\sk{प्रकटयति}}, \emph{prakaṭayati}.) Guilt? (\textbf{\sk{दोषः}}, \emph{doṣaḥ}.) Punishes? (\textbf{\sk{दण्डयति}}, \emph{daṇḍayati}.) Sets free? (\textbf{\sk{मोचयति}}, \emph{mocayati}.) A prison? (\textbf{\sk{कारागारम्}}, \emph{kārāgāram}.) A chain? (\textbf{\sk{शृंखला}}, \emph{śṛṅkhalā}.)
\end{tcolorbox}
